% ---
% RESUMOS
% ---

% resumo em português
\setlength{\absparsep}{18pt} % ajusta o espaçamento dos parágrafos do resumo

\begin{resumo}
Este trabalho investiga a deteccao de intrusoes em redes CAN de sistemas veiculares no contexto da Internet of Vehicles (IoVT). O objetivo foi comparar quatro modelos supervisionados de aprendizado de maquina para classificacao multiclasse de trafego CAN: Regressao Logistica, SVM, MLP e XGBoost. O estudo utilizou um conjunto final alinhado e limpo com cinco classes (BENIGN, DoS, Fuzzy, GEAR e RPM), adotando como base canonica os artefatos finais de avaliacao em \texttt{results/metrics/*\_final\_clean}.

As metricas principais de comparacao foram accuracy, precision macro, recall macro e F1 macro, complementadas por analise por classe. Os resultados mostraram desempenho global superior do XGBoost (accuracy 0.99998; F1 macro 0.99998), seguido do MLP (accuracy 0.99908; F1 macro 0.99905). Regressao Logistica e SVM apresentaram desempenho inferior aos dois melhores modelos, embora com resultados consistentes no cenario avaliado. A classe Fuzzy foi a mais desafiadora para todos os classificadores, enquanto GEAR e RPM apresentaram separacao quase perfeita.

Como contribuicao, o trabalho oferece uma comparacao reproduzivel entre modelos para IDS em CAN, com rastreabilidade dos resultados e alinhamento entre metodologia e avaliacao final.


 \textbf{Palavras-chave}: IDS, rede CAN, IoVT, aprendizado de maquina, deteccao de intrusoes.
\end{resumo}

% resumo em inglês
\begin{resumo}[Abstract]
 \begin{otherlanguage*}{english}
This work investigates intrusion detection in in-vehicle CAN networks within the Internet of Vehicles (IoVT) context. The main objective was to compare four supervised machine learning models for multiclass CAN traffic classification: Logistic Regression, SVM, MLP, and XGBoost. The study used a final aligned and cleaned dataset with five classes (BENIGN, DoS, Fuzzy, GEAR, and RPM), adopting the final evaluation artifacts in \texttt{results/metrics/*\_final\_clean} as the canonical source.

The main comparison metrics were accuracy, macro precision, macro recall, and macro F1, supported by class-level analysis. Results show that XGBoost achieved the best overall performance (accuracy 0.99998; macro F1 0.99998), followed by MLP (accuracy 0.99908; macro F1 0.99905). Logistic Regression and SVM performed below the two top models, although both remained consistent in the evaluated scenario. The Fuzzy class was the most challenging class across all classifiers, while GEAR and RPM were almost perfectly separable.

As a contribution, this work provides a reproducible CAN IDS model comparison with traceable results and strong alignment between methodology and final evaluation.
   \vspace{\onelineskip}
 
   \noindent 
  \textbf{Keywords}: IDS, CAN network, IoVT, machine learning, intrusion detection.
 \end{otherlanguage*}
\end{resumo}

%% resumo em francês 
%\begin{resumo}[Résumé]
% \begin{otherlanguage*}{french}
%    Il s'agit d'un résumé en français.
% 
%   \textbf{Mots-clés}: latex. abntex. publication de textes.
% \end{otherlanguage*}
%\end{resumo}
%
%% resumo em espanhol
%\begin{resumo}[Resumen]
% \begin{otherlanguage*}{spanish}
%   Este es el resumen en español.
%  
%   \textbf{Palabras clave}: latex. abntex. publicación de textos.
% \end{otherlanguage*}
%\end{resumo}
% ---
